\makeatletter
\def\input@path{{./}{../}{../../}{../../../}{../../../../}{preamble/}{../preamble/}{../../preamble/}{../../../preamble/}{../../../../book/preamble/}}
\makeatother
\input{preamble_exams.tex}
\input{graphs/highlights.tex}
% \GTenHighlightsfalse % Uncomment to disable instructional highlights.
% Population/model, sample evidence, design threshold, and result roles.
\newcommand{\Pop}[1]{\HighlightStroke{GTenHLParameter}{#1}}
\newcommand{\Obs}[1]{\HighlightStroke{GTenHLValue}{#1}}
\newcommand{\Design}[1]{\HighlightStroke{GTenHLThird}{#1}}
\newcommand{\Result}[1]{\HighlightStroke{GTenHLOutcome}{#1}}

\title{Grade 11 Global Economics\\Term 3 --- Lesson 1\\Home Practice C1: Identifying the ingredients of a test}
\author{}\date{}
\hypersetup{pdftitle={Grade 11 Global Economics - Term 3 - Lesson 1 - C1},pdfsubject={Identifying the ingredients of a test}}
\begin{document}
\maketitle
\setcounter{tocdepth}{1}
\setcounter{secnumdepth}{2}
\phantomsection\label{toc}
\tableofcontents
\section*{Evaluated criteria}
\begin{description}
\item[\textbf{C1}] Identifies the essential elements required to conduct a hypothesis test in finance and economics, including the population parameter of interest, hypothesised value, sample statistic, sample size, variability, significance level, and relevant sampling conditions.
\end{description}
\textbf{How to practise.} Build an information inventory; separate population quantities from observed summaries. Attempt each task before opening its solution link. Use the return arrows to revisit the problem, group, or contents. All datasets and scenarios in this activity are illustrative.
\clearpage
\section{Introduction}\label{sec:introduction}
\subsection*{Highlight legend}
\phantomsection\label{tab:highlight-legend}
\begin{tabularx}{\linewidth}{@{}>{\raggedright\arraybackslash}p{.22\linewidth}>{\raggedright\arraybackslash}p{.25\linewidth}>{\raggedright\arraybackslash}X@{}}
\toprule
\textbf{Symbol or notation}&\textbf{Concept}&\textbf{Brief meaning}\\\midrule
$\Pop{\mu},\Pop{p}$ & \Pop{Parameter} & Population feature. \\
$\Obs{\bar x},\Obs{\hat p}$ & \Obs{Statistic} & Sample evidence. \\
$\Pop{\mu_0},\Pop{p_0}$ & \Pop{Benchmark} & Reference value. \\
$\Obs{n},\Pop{\sigma},\Obs{s}$ & \Obs{Size and spread} & Amount and variability. \\
$\Design{\alpha}$ & \Design{Level} & Planned threshold. \\
\bottomrule
\end{tabularx}
\subsection{Introductory example 1: Invoice values}\label{ex:intro-one}
A random independent sample of 36 invoices has $\Obs{\bar x}=\Obs{79}$ USD. The reference is $\Pop{\mu_0}=\Pop{75}$ USD; known population $\Pop{\sigma}=\Pop{12}$ USD. The planned level is $\Design{\alpha}=\Design{0.05}$. Unknown $\mu$ describes all eligible invoices.
\subsection{Introductory example 2: Loan defaults}\label{ex:intro-two}
Among 200 independently sampled loans, 12 default within one year, so $\hat p=12/200=0.06$. The reference is $p_0=0.05$ and the level is 0.01. Observed 6\% does not establish the population rate.
\section{Notation Reference and Practice Procedure}\label{sec:notation-and-procedure}
\subsection{Notation reference}\label{tab:notation-reference}
\begin{tabularx}{\linewidth}{@{}>{\raggedright\arraybackslash}p{.16\linewidth}>{\raggedright\arraybackslash}p{.18\linewidth}>{\raggedright\arraybackslash}p{.28\linewidth}>{\raggedright\arraybackslash}X@{}}
\toprule
\textbf{Symbol or notation}&\textbf{Concept}&\textbf{Definition}&\textbf{Meaning in the introductory examples}\\\midrule
$\Pop{\mu},\Pop{p}$ & \Pop{Parameter} & $\mu$ is a population mean; $p$ a population proportion. & \textbullet\ \hyperref[ex:intro-one]{Example 1}: mean invoice value.\par\textbullet\ \hyperref[ex:intro-two]{Example 2}: default proportion. \\
$\Obs{\bar x},\Obs{\hat p}$ & \Obs{Statistic} & $\bar x=\sum x_i/n$; $\hat p=x/n$ for $x$ successes. & \textbullet\ \hyperref[ex:intro-one]{Example 1}: 79 USD.\par\textbullet\ \hyperref[ex:intro-two]{Example 2}: $12/200=0.06$. \\
$\Pop{\mu_0},\Pop{p_0}$ & \Pop{Benchmark} & Hypothesised population value, separate from the sample estimate. & \textbullet\ \hyperref[ex:intro-one]{Example 1}: 75 USD.\par\textbullet\ \hyperref[ex:intro-two]{Example 2}: 0.05. \\
$\Obs{n},\Pop{\sigma},\Obs{s}$ & \Obs{Size and spread} & $n$ counts units; $\sigma$ is population SD; $s$ estimates it. Variance is SD squared. & \textbullet\ \hyperref[ex:intro-one]{Example 1}: $n=36$, $\sigma=12$ USD.\par\textbullet\ \hyperref[ex:intro-two]{Example 2}: $n=200$. \\
$\Design{\alpha}$ & \Design{Level} & Chosen before results; controls false rejection under the null model. & \textbullet\ \hyperref[ex:intro-one]{Example 1}: 0.05.\par\textbullet\ \hyperref[ex:intro-two]{Example 2}: 0.01. \\
\bottomrule
\end{tabularx}
\subsection{Practice procedure}
\begin{enumerate}
\item Define the population, unit, measurement, and period.
\item Name the parameter and benchmark with units.
\item Record statistic, sample size, and whether spread is known or estimated.
\item Record level, random sampling, independence, and distribution conditions.
\item State missing information; a full test is not required.
\end{enumerate}
\section{General Worked Problems}\label{sec:general-worked}
\SectionProblemsTableOfContents{\SectionProblemLink{prob:inventory}{Payment-platform inventory} \\
\SectionProblemLink{prob:spread}{Known versus estimated spread}}
\subsection{Payment-platform inventory}\label{prob:inventory}
A provider claims 62\% of firms use its platform weekly. A simple random sample of 200 distinct firms from 20,000 eligible firms finds 112 users. The level is 5\%. Firms respond independently and once each.
\QuestionsTasks[sec:general-worked]{{1}/{sol:inventory}/{Identify the test ingredients and distinguish population quantities from sample summaries.}}
\subsubsection*{Worked solution}\label{sol:inventory}
The unit is a firm. $\Pop{p}$ is the unknown population weekly-use proportion. $\Pop{p_0}=\Pop{0.62}$, $x=112$, $\Obs{n}=\Obs{200}$, and $\Obs{\hat p}=112/200=\Obs{0.56}$. Both $x$ and $\hat p$ describe the sample only. $\Design{\alpha}=\Design{0.05}$. Null variability of $\hat p$ is $\sqrt{p_0(1-p_0)/n}$. Random sampling is stated; $200<0.10(20{,}000)$ supports approximate independence without replacement. Expected counts $200(0.62)=124$ and $200(0.38)=76$ both exceed 10. This inventory supports model selection, but is not a completed test.
\secInicio[prob:inventory][sec:general-worked]
\subsection{Known versus estimated spread}\label{prob:spread}
A random independent sample of 36 export invoices gives $\bar x=79$ USD against 75 USD. Report A supplies known population variance 144 USD$^2$; report B supplies only sample variance 144 USD$^2$. The population is approximately normal; $\alpha=0.01$.
\QuestionsTasks[sec:general-worked]{{1}/{sol:spread}/{Create both inventories and explain the difference without testing.}}
\subsubsection*{Worked solution}\label{sol:spread}
Both concern $\Pop{\mu}$, population mean invoice value in USD: $\mu_0=75$, $\Obs{\bar x}=\Obs{79}$, $n=36$, and $\alpha=0.01$. In A, $\sigma^2=144$ gives known $\Pop{\sigma}=\Pop{12}$ USD. In B, $s^2=144$ gives estimated $\Obs{s}=\Obs{12}$ USD; population variance remains unknown. Equal numbers do not give equal information. Random independent sampling and approximate normality are supplied. Both use a mean model, but B retains uncertainty from estimating variability.
\secInicio[prob:spread][sec:general-worked]
\section{Advanced Problems}\label{sec:advanced}
\SectionProblemsTableOfContents{\SectionProblemLink{prob:missing}{An incomplete investment report}}
\subsection{An incomplete investment report}\label{prob:missing}
A report says investment success is 58\% against a 50\% benchmark and concludes improvement. No count, sampling plan, success definition, or level is supplied. Another paragraph gives average profit 40 USD without spread.
\QuestionsTasks[sec:advanced]{{1}/{sol:missing}/{Draft an information request for both claims and explain what is currently supported.}}
\subsubsection*{Worked solution}\label{sol:missing}
For success, define the population, horizon, and binary success rule; obtain exact $x$ and $n$, confirm that 58\% is $\hat p$ and that 50\% was fixed independently. Request the sampling method, dependence information, and planned $\alpha$. These allow independence and expected-count conditions to be assessed. For profit, define $\mu$ and its benchmark; request $n$, confirmation that 40 USD is $\bar x$, and known $\sigma$ or sample $s$, plus population shape and outliers. A success proportion cannot supply dollar-profit variability. These descriptive statements alone do not establish population improvement; missing quantities must not be guessed.
\secInicio[prob:missing][sec:advanced]
\end{document}

